% generated by the concordance script (main thread); do not hand-edit
--- & Tier 3 --- Astrophysics and the Verdict & 1 & title page, \hyperref[sec:covers]{What this tier covers} \\
--- & Status of this document & 47 & \hyperref[sec:status]{Status of this document} \\
--- & How to study each node & 122 & \hyperref[sec:howtostudy]{How to study each node} \\
--- & Contents & 148 & \S\ref{app:sectionmap} \\
--- & Part 0 --- The astrophysics this project is for & 228 & \S\ref{part:engine}--\S\ref{part:host} \\
0.1 & The star as a self-gravitating thermodynamic engine & 256 & \S\ref{sec:engine} \\
0.1.1 & Hydrostatic equilibrium, and the pressure a star needs & 262 & \S\ref{sec:hydrostatic} \\
0.1.2 & The virial theorem, and the negative heat capacity --- derived & 296 & \S\ref{sec:virial} \\
0.1.3 & What temperature a contracting core reaches & 351 & \S\ref{sec:tcore} \\
0.1.4 & Degeneracy: the wall the star runs into & 398 & \S\ref{sec:degeneracy} \\
0.2 & The burning staircase, derived rather than tabulated & 443 & \S\ref{sec:staircase} \\
0.2.1 & Why each fuel ignites where it does & 450 & \S\ref{sec:ignition} \\
0.2.2 & The energy per gram of each stage, computed & 574 & \S\ref{sec:energypergram} \\
0.2.3 & The neutrino clock --- why the last stages are short & 612 & \S\ref{sec:neutrinoclock} \\
0.2.4 & What the star is doing while the network is being called & 711 & \S\ref{sec:whilecalled} \\
0.2.5 & Convection --- the transport a one-zone burner cannot see & 734 & \S\ref{sec:convection} \\
0.3 & Silicon burning proper & 871 & \S\ref{part:siburn} \\
0.3.1 & The mechanism, mechanically & 878 & \S\ref{sec:mechanism} \\
0.3.2 & The two clusters, and what the group boundary means physically & 942 & \S\ref{sec:twoclusters} \\
0.3.3 & The endpoint, and why \Ye{} chooses it & 1027 & \S\ref{sec:endpoint} \\
0.3.4 & Hydrostatic versus explosive silicon burning & 1067 & \S\ref{sec:hydroexplosive} \\
0.3.5 & $\alpha$-rich freeze-out --- the third regime & 1124 & \S\ref{sec:alpharich} \\
0.3.6 & The NSE transition --- where a network is replaced by a table & 1169 & \S\ref{sec:nsetransition} \\
0.4 & The onion, the iron core, and the shell that builds it & 1261 & \S\ref{sec:ironcore} \\
0.4.1 & The structure at the end & 1263 & \S\ref{sec:onion} \\
0.4.2 & The effective Chandrasekhar mass & 1309 & \S\ref{sec:mch} \\
0.4.3 & Compactness: the structural predictor & 1373 & \S\ref{sec:compactness} \\
0.4.4 & What pre-supernova \Ye{} actually looks like & 1407 & \S\ref{sec:presnye} \\
0.4.5 & Which stars are these, actually? & 1444 & \S\ref{sec:whichstars} \\
0.5 & Core collapse, derived end to end & 1499 & \S\ref{sec:corecollapse} \\
0.5.1 & The instability: $\Gamma$$_1$ < 4/3 & 1505 & \S\ref{sec:gamma43} \\
0.5.2 & The two runaways & 1539 & \S\ref{sec:runaways} \\
0.5.3 & The free-fall timescale & 1609 & \S\ref{sec:freefall} \\
0.5.4 & Neutrino trapping --- derived, with numbers & 1636 & \S\ref{sec:trapping} \\
0.5.5 & The homologous core and the bounce & 1700 & \S\ref{sec:homologous} \\
0.5.6 & Why the prompt shock fails --- the energy budget & 1734 & \S\ref{sec:promptshock} \\
0.5.7 & The delayed neutrino-driven mechanism & 1794 & \S\ref{sec:delayed} \\
0.5.8 & The mass cut, and the honest caveat & 1842 & \S\ref{sec:masscut} \\
0.5.9 & What "explodability" actually means operationally & 1876 & \S\ref{sec:explodability} \\
0.6 & The observable end & 1937 & \S\ref{part:observables} \\
0.6.1 & \nuc{Ni}{56} and the light curve & 1943 & \S\ref{sec:ni56} \\
0.6.2 & Isotopic ratios and the neutron excess & 1997 & \S\ref{sec:isotopicratios} \\
0.6.3 & Galactic chemical evolution & 2047 & \S\ref{sec:gce} \\
0.6.4 & Remnants: the two objects everyone cites & 2070 & \S\ref{sec:remnants} \\
0.6.5 & \warnglyph{} The derivative nobody has assembled & 2095 & \S\ref{sec:missingderivative} \\
0.6.6 & The error hierarchy --- what the emulator is actually competing with & 2153 & \S\ref{sec:errorhierarchy} \\
\textbf{0.7} & \textbf{Where \Ye{} comes from --- the weak sector, quantitatively} & 2213 & \S\ref{part:weak} \\
\textbf{0.7.1} & \textbf{The identity, and what can move it} & 2218 & \S\ref{sec:yeidentity} \\
\textbf{0.7.2} & \textbf{Electron capture in a degenerate plasma --- why it is (nearly) a one-way ratchet} & 2267 & \S\ref{sec:ecratchet} \\
\textbf{0.7.3} & \textbf{The measured \Ye{} carriers} & 2322 & \S\ref{sec:yecarriers} \\
0.8 & The computational setting & 2353 & \S\ref{part:host} \\
0.8.1 & What MESA actually does per timestep & 2357 & \S\ref{sec:mesastep} \\
0.8.2 & The cost model, and the compromise everybody makes & 2403 & \S\ref{sec:costmodel} \\
0.8.3 & The emulation landscape & 2447 & \S\ref{sec:landscape} \\
0.8.4 & What this project adds, stated narrowly & 2476 & \S\ref{sec:whatadds} \\
0.8.5 & \warnglyph{} The bottleneck premise has never been measured & 2502 & \S\ref{sec:bottleneck} \\
0.8.6 & \warnglyph{} The call-site contract --- seven requirements nobody has written down & 2554 & \S\ref{sec:callsite} \\
0.9 & Astrophysics $\to$ gates: the mapping table & 2664 & \S\ref{sec:gatemap} \\
\textbf{0.10} & \textbf{Self-check for Part 0} & 2714 & Self-checks of \S\ref{part:engine}, \S\ref{part:siburn}, \S\ref{part:weak}, \S\ref{part:collapse}, \S\ref{part:host} (App.~\ref{app:selfcheckmoves}) \\
--- & Part I (S11) --- bbq and MESA as operated & 2758 & \S\ref{part:generator} \\
I.1 & What a one-zone burner is --- and what it deliberately is not & 2765 & \S\ref{sec:onezone} \\
\textbf{I.2} & \textbf{MESA's network solver, in enough detail to read the config} & 2799 & \S\ref{sec:netsolver} \\
\textbf{I.2b} & \textbf{What the network returns --- the interface, as the host sees it} & 2843 & \S\ref{sec:netreturns} \\
I.3 & bbq's four modes, and why only one made the shipped trajectories & 2907 & \S\ref{sec:bbqmodes} \\
I.4 & The inlist as a configuration surface & 2939 & \S\ref{sec:inlist} \\
I.5 & The shipped-label pipeline, reconstructed & 2988 & \S\ref{sec:shippedpipeline} \\
I.6 & Why a rerun campaign was necessary & 3028 & \S\ref{sec:campaign} \\
I.6b & The campaign's coverage geometry --- and what 209 runs can and cannot say & 3083 & \S\ref{sec:coverage} \\
I.7 & Execution & 3132 & \S\ref{sec:execution} \\
I.8 & Validation --- four checks, and what each was for & 3152 & \S\ref{sec:validation} \\
I.9 & Self-check for S11 & 3202 & Self-checks of \S\ref{part:engine}, \S\ref{part:siburn}, \S\ref{part:weak}, \S\ref{part:collapse}, \S\ref{part:host} (App.~\ref{app:selfcheckmoves}) \\
--- & Part II (S12) --- Trajectories and the eps\_nuc pin & 3219 & \S\ref{part:trajectories} \\
II.1 & The file format & 3225 & \S\ref{sec:fileformat} \\
II.1b & What a trajectory row is \emph{not} & 3250 & \S\ref{sec:rowisnot} \\
II.2 & The composition route to energy --- derived & 3283 & \S\ref{sec:compositionroute} \\
II.3 & The pin, as a hypothesis test & 3327 & \S\ref{sec:pin} \\
II.4 & The stall, and its two reinterpretations & 3390 & \S\ref{sec:stall} \\
II.5 & Self-check for S12 & 3445 & Self-checks of \S\ref{part:engine}, \S\ref{part:siburn}, \S\ref{part:weak}, \S\ref{part:collapse}, \S\ref{part:host} (App.~\ref{app:selfcheckmoves}) \\
--- & Part III (S13) --- Silicon burning on real histories, and the label pathology & 3457 & \S\ref{part:pathology} \\
III.1 & What a real trajectory looks like & 3463 & \S\ref{sec:realtrajectory} \\
III.2 & The census: how far are the labels from NSE? & 3498 & \S\ref{sec:census} \\
III.3 & The mechanism: gh-575, traced from a Fortran branch & 3538 & \S\ref{sec:gh575} \\
III.4 & The witness experiment & 3605 & \S\ref{sec:witness} \\
III.4b & How this could have been caught on day one & 3644 & \S\ref{sec:dayone} \\
III.5 & The consequence: \Ye{} itself is wrong & 3689 & \S\ref{sec:yewrong} \\
III.6 & The epistemics: what a negative result about your own dataset obliges & 3719 & \S\ref{sec:epistemics} \\
III.7 & What the pathology does \emph{not} invalidate & 3795 & \S\ref{sec:notinvalidated} \\
III.8 & Self-check for S13 & 3822 & Self-checks of \S\ref{part:engine}, \S\ref{part:siburn}, \S\ref{part:weak}, \S\ref{part:collapse}, \S\ref{part:host} (App.~\ref{app:selfcheckmoves}) \\
--- & Part IV (S14) --- The kill-test & 3841 & \S\ref{part:killtest} \\
IV.1 & The question, restated & 3846 & \S\ref{sec:question} \\
IV.2 & The instruments & 3874 & \S\ref{sec:instruments} \\
IV.3 & The two distributions --- and why they are different worlds & 3929 & \S\ref{sec:twodistributions} \\
IV.4 & Gate 1 --- \Ye{} coverage: a structural pass & 3981 & \S\ref{sec:gate1} \\
IV.5 & Gate 5 --- the Guidry sweep returns EMPTY & 4007 & \S\ref{sec:gate5} \\
IV.6 & Gate 2 --- the split verdict, and the 1/$\kappa$ amplification & 4062 & \S\ref{sec:gate2} \\
IV.7 & Gates 3, 7, 8 --- spread, timescales, energy & 4106 & \S\ref{sec:gates378} \\
IV.8 & The verdict, and the statistics it does not carry & 4158 & \S\ref{sec:verdict} \\
IV.8b & What a complete statistics package would look like & 4198 & \S\ref{sec:statspackage} \\
IV.9 & What would have changed the answer & 4262 & \S\ref{sec:whatwouldflip} \\
IV.10 & Self-check for S14 & 4281 & Self-checks of \S\ref{part:engine}, \S\ref{part:siburn}, \S\ref{part:weak}, \S\ref{part:collapse}, \S\ref{part:host} (App.~\ref{app:selfcheckmoves}) \\
--- & Part V (S15) --- The ML target, decided & 4300 & \S\ref{part:target} \\
V.1 & The two targets, restated & 4305 & \S\ref{sec:twotargets} \\
V.2 & What the verdict decided, and what it left open & 4347 & \S\ref{sec:decided} \\
V.3 & Why conservation must live in the decode step --- the NuGNN failure mode & 4375 & \S\ref{sec:decodestep} \\
V.4 & $\Phi$ supervision: proven where it matters, blocked where it does not & 4421 & \S\ref{sec:phisupervision} \\
V.5 & The identifiability cost nobody has paid attention to & 4463 & \S\ref{sec:identifiability} \\
V.5b & The output space is a simplex, and nothing in the project treats it as one & 4491 & \S\ref{sec:simplex} \\
V.5c & Gradients through the conservation map --- two different questions & 4559 & \S\ref{sec:gradients} \\
V.6 & Self-check for S15 & 4595 & Self-checks of \S\ref{part:engine}, \S\ref{part:siburn}, \S\ref{part:weak}, \S\ref{part:collapse}, \S\ref{part:host} (App.~\ref{app:selfcheckmoves}) \\
--- & Part VI (S16) --- The baseline you are trying to beat & 4612 & \S\ref{part:baseline} \\
VI.1 & What the NNN actually is & 4617 & \S\ref{sec:nnnis} \\
VI.2 & Reading that design critically & 4674 & \S\ref{sec:nnncritical} \\
VI.3 & What the NNN achieves & 4729 & \S\ref{sec:nnnachieves} \\
VI.4 & NuGNN, the closer competitor & 4758 & \S\ref{sec:nugnn} \\
VI.5 & What "beating it" has to mean & 4791 & \S\ref{sec:beating} \\
VI.5b & The protocol, written out & 4838 & \S\ref{sec:protocol} \\
VI.6 & Self-check for S16 & 4895 & Self-checks of \S\ref{part:engine}, \S\ref{part:siburn}, \S\ref{part:weak}, \S\ref{part:collapse}, \S\ref{part:host} (App.~\ref{app:selfcheckmoves}) \\
--- & Part VII --- What Tier 3 hands forward & 4911 & \S\ref{part:handoff} \\
VII.1 & What Phase 0 established & 4913 & \S\ref{sec:established} \\
VII.2 & What Phase 0 retired & 4937 & \S\ref{sec:retired} \\
VII.3 & The four escalations & 4955 & \S\ref{sec:escalations} \\
VII.4 & The state Tier 4 inherits & 4971 & \S\ref{sec:inherits} \\
VII.5 & The three ideas to carry, if you carry nothing else & 5004 & \S\ref{sec:threeideas} \\
--- & Part VIII --- The open register & 5028 & \S\ref{part:register} \\
VIII.1 & The register & 5036 & \S\ref{sec:register} \\
--- & The second-pass items (Part IX) & 5059 & \S\ref{sec:register-second} \\
--- & The third-pass items (Part X) --- the emulator as an object & 5078 & \S\ref{sec:register-third} \\
VIII.2 & Ranked --- the top eight & 5098 & \S\ref{sec:ranked} \\
VIII.3 & How this register was built & 5142 & \S\ref{sec:registerbuilt} \\
\textbf{---} & \textbf{Part IX --- Prerequisites outside the map: a cross-tier audit} & 5204 & \S\ref{sec:boundaryaudit} \\
IX.1 & What was swept, and how & 5212 & \S\ref{sec:swept} \\
IX.2 & The areas with no study node anywhere & 5237 & \S\ref{sec:noareas} \\
IX.3 & Three developed here, having no home section & 5258 & \S\ref{sec:threedeveloped} \\
IX.3.1 & Experiment design and the multiplicity problem & 5260 & \S\ref{sec:multiplicity} \\
IX.3.2 & The artifact lifecycle --- what actually ships & 5299 & \S\ref{sec:artifact} \\
IX.3.3 & Nuclear-data uncertainty is probably the binding term --- promote it & 5338 & \S\ref{sec:nucleardata} \\
IX.4 & Axes checked and found already covered & 5373 & \S\ref{sec:axescovered} \\
IX.5 & Is the audit converging? & 5402 & \S\ref{sec:converging} \\
\textbf{---} & \textbf{Part X --- The third pass: the emulator as an object} & 5429 & \S\ref{sec:emulatorobject} \\
X.1 & The emulator has its own fixed points, and nobody has asked what they are & 5441 & \S\ref{sec:fixedpoints} \\
X.2 & The learned Jacobian --- does stiffness survive learning? & 5521 & \S\ref{sec:jacobian} \\
X.3 & The fallback gate is a discontinuity in the host's implicit solve & 5563 & \S\ref{sec:dispatcher} \\
X.4 & $\rho$ is not a free input dimension & 5605 & \S\ref{sec:rhoinput} \\
X.5 & What ``calibrated'' means for a deterministic map & 5654 & \S\ref{sec:calibrated} \\
X.6 & Rates as inputs --- the capability nobody claimed & 5711 & \S\ref{sec:ratesinputs} \\
X.7 & Which $\Delta$t does the host actually use? & 5746 & \S\ref{sec:hostdt} \\
X.8 & Sample complexity --- a million points in 82 dimensions & 5768 & \S\ref{sec:samplecomplexity} \\
X.9 & Two smaller ones & 5803 & \S\ref{sec:twosmaller} \\
X.10 & One unretired premise, found while sweeping & 5823 & \S\ref{sec:unretired} \\
X.11 & Is the audit converging \emph{now}? & 5867 & \S\ref{sec:convergingnow} \\
--- & Where Tier 3 hands off & 5915 & \S\ref{sec:handsoff} \\
--- & References & 5933 & \hyperref[sec:references]{References} \\
